% Sección: derivación del Cross-Entropy Method.
\section{Optimización: Cross-Entropy Method}
\label{sec:cem}

Optimizamos $\max_{w} J(w)$ con una familia de muestreo Gaussiana
$g(w; \theta) = \mathcal{N}(w; \mu, \Sigma)$, con $\theta = (\mu, \Sigma)$.

\subsection{Densidad objetivo y proyección por entropía cruzada}
Para un nivel de desempeño $\lambda$, la densidad ideal concentra masa en la
región élite:
\[
  q_{\lambda}(w) \propto \mathbf{1}\{J(w) \ge \lambda\}\, g(w; \theta).
\]
Buscamos $\theta$ que minimice la divergencia de Kullback--Leibler (entropía
cruzada) hacia $q_{\lambda}$:
\[
  \theta^{*}
  = \arg\min_{\theta} D_{\mathrm{KL}}\!\big(q_{\lambda} \,\|\, g(\cdot; \theta)\big)
  = \arg\max_{\theta} \mathbb{E}_{q_{\lambda}}\!\left[\log g(w; \theta)\right].
\]
Con muestras $w_{1}, \dots, w_{N} \sim g(\cdot; \theta_{t})$ y nivel adaptativo
$\lambda_{t}$ igual al cuantil $(1-\rho)$ de $\{J(w_{i})\}$, el conjunto élite es
$\mathcal{E}_{t} = \{w_{i} : J(w_{i}) \ge \lambda_{t}\}$, con
$|\mathcal{E}_{t}| = \lceil \rho N \rceil$.

\subsection{Solución en forma cerrada}
El paso de maximización es un ajuste Gaussiano por máxima verosimilitud sobre
los élites. Partiendo de
$\log \mathcal{N}(w; \mu, \Sigma)
 = -\tfrac{1}{2}\log|2\pi\Sigma| - \tfrac{1}{2}(w-\mu)^{\top}\Sigma^{-1}(w-\mu)$
y anulando las derivadas,
\begin{align}
  \frac{\partial}{\partial \mu} \sum_{w \in \mathcal{E}} \log \mathcal{N}
    &= \Sigma^{-1} \sum_{w \in \mathcal{E}} (w - \mu) = 0
    &&\Rightarrow&
    \mu^{*} &= \frac{1}{|\mathcal{E}|} \sum_{w \in \mathcal{E}} w, \\
  \frac{\partial}{\partial \Sigma^{-1}} \sum_{w \in \mathcal{E}} \log \mathcal{N}
    &= \frac{|\mathcal{E}|}{2}\Sigma
       - \frac{1}{2}\sum_{w \in \mathcal{E}}(w-\mu^{*})(w-\mu^{*})^{\top} = 0
    &&\Rightarrow&
    \Sigma^{*} &= \frac{1}{|\mathcal{E}|}
      \sum_{w \in \mathcal{E}} (w-\mu^{*})(w-\mu^{*})^{\top}.
\end{align}
La actualización que minimiza la entropía cruzada es, pues, la media y la
covarianza empíricas del conjunto élite. En la práctica se usa
$\Sigma = \mathrm{diag}(\sigma_{1}^{2}, \dots, \sigma_{d}^{2})$ por eficiencia.

\subsection{Inyección de ruido y evaluación Monte Carlo}
El CEM puro hace $\Sigma_{t} \to 0$ y colapsa a un óptimo local; para mantener
exploración se añade ruido decreciente:
\[
  \Sigma_{t+1} \leftarrow \Sigma^{*} + Z_{t} I,
  \qquad Z_{t} = \max\!\big(0,\; a - t/b\big).
\]
Cada $J(w_{i})$ se estima por Monte Carlo sobre $m$ episodios,
$\hat{J}(w_{i}) = \tfrac{1}{m} \sum_{k=1}^{m} G^{(k)}(w_{i})$, con
$\mathrm{Var}[\hat{J}] \propto 1/m$ \citep{rubinstein1999, szita2006}.

\subsection{Implementación}
En este trabajo se emplea covarianza diagonal y el ruido decrece de forma lineal,
$Z_{t} = c\,(1 - t/T)$, con $T$ el número de generaciones. La evaluación de cada
candidato usa un \emph{tope de colocaciones} por partida (una buena política juega
partidas casi ilimitadas) y \emph{números aleatorios comunes}: todos los candidatos
de una misma generación se evalúan con las mismas semillas de partida, lo que
reduce la varianza de la comparación. El vector final se normaliza a
$\lVert w \rVert = 1$; al ser la política invariante a escala, esto no altera el
desempeño pero mantiene los pesos acotados, en contraste con la divergencia de
magnitud observada en TD (Sección~\ref{sec:td}).
